\section{课程背景与回顾}

本讲是 KAIST CS492D（扩散模型与生成模型）的第 9 讲，主题为 \textbf{Flow Matching}（流匹配）。在前几讲中，我们已经详细讨论了扩散模型（Diffusion Models）的各种变体——从 DDPM 到 DDIM，再到基于 SDE/ODE 的 Score-Based 模型和 DPM-Solver 等加速技术。本讲将从一个全新的角度引入生成模型：\textbf{Continuous Normalizing Flows}（连续正则化流）和 \textbf{Flow Matching}（流匹配）。

\begin{knowledgebox}{课程脉络回顾}
在前面的课程中，我们建立了以下认知链条：
\begin{enumerate}
    \item \textbf{DDPM}：离散时间步的去噪扩散模型，典型需要 1000 步采样
    \item \textbf{DDIM}：确定性采样，将步数降至 50--200 步
    \item \textbf{Score-Based SDE}：将扩散过程统一为随机微分方程（SDE）框架
    \item \textbf{Probability Flow ODE}：从 SDE 导出等价的 ODE 形式，实现确定性生成
    \item \textbf{DPM-Solver}：利用 ODE 结构的专用求解器，进一步将步数降至 10--15 步
\end{enumerate}
本讲的 Flow Matching 将提供另一种理解和训练 ODE-based 生成模型的方式。
\end{knowledgebox}

\subsection{扩散模型的速度瓶颈}

扩散模型的核心优势在于能够生成高质量、无模式崩塌（mode collapse）的输出。但其主要瓶颈始终是\textbf{采样速度}：

\begin{table}[H]
\centering
\begin{tabular}{lcc}
\toprule
\textbf{方法} & \textbf{典型步数} & \textbf{是否需要重新训练} \\
\midrule
DDPM & $\sim$1000 & 否 \\
DDIM & 50--200 & 否 \\
DPM-Solver (高阶) & 10--15 & 否 \\
Consistency Models & 1--2 & 是（微调） \\
Distillation 方法 & 1--5 & 是（微调） \\
\bottomrule
\end{tabular}
\caption{扩散模型采样加速方法对比}
\end{table}

讲者指出，DDPM 到 DPM-Solver 的演进实现了约 \textbf{100 倍}的加速，且不需要重新训练模型——这是扩散模型非常有吸引力的特性。但要进一步将步数降至 1--5 步，则需要 distillation 等额外训练手段。

\subsection{Consistency Models 简介}

\begin{importantbox}{Consistency Models 的核心思想}
Consistency Models 的基本思想是：对于 ODE 轨迹上的所有点，训练网络使其能从轨迹的\textbf{任意中间点}直接预测最终输出。具体来说：
\begin{enumerate}
    \item 从标准高斯样本出发，运行 ODE 求解器获得完整轨迹
    \item 对轨迹上的每个中间点 $x_t$，利用 Tweedie 公式计算最终输出的期望
    \item 训练网络使得较早时刻的预测与较晚时刻的预测保持一致（consistent）
\end{enumerate}
这样，在推理时只需一步或两步即可从噪声直接跳到数据样本。
\end{importantbox}

Consistency Models 和其他 distillation 方法的共同缺点是：它们都需要预训练的扩散模型作为起点，通过额外的微调来实现加速。这意味着需要在训练阶段投入更多计算资源，以换取测试时更快的速度。

\subsection{本章小结}

从 DDPM 到 Consistency Models，扩散模型社区在采样加速方面取得了巨大进展。然而，这些方法要么受限于特定的噪声调度（noise schedule），要么需要复杂的蒸馏训练。Flow Matching 将从一个完全不同的角度出发，提供更灵活、更简洁的生成模型训练框架。

